\documentclass[11pt,letterpaper]{article}
\usepackage[margin=0.78in,headheight=15pt,headsep=17pt]{geometry}
\usepackage[T1]{fontenc}
\usepackage{lmodern,amsmath,amssymb,mathtools,microtype}
\usepackage{booktabs,tabularx,array,enumitem,xcolor,fancyhdr}
\definecolor{inkblue}{RGB}{27,56,82}
\usepackage[colorlinks=true,allcolors=inkblue]{hyperref}
\hypersetup{pdftitle={Springer--Borel maps: a structural research program},
pdfauthor={Research program for Sean Rogers},
pdfsubject={Integral sources, characteristic images, generic degrees, equivariant and generalized cohomology}}
\pagestyle{fancy}\fancyhf{}
\fancyhead[L]{\small Springer--Borel maps: research proposal}
\fancyhead[R]{\small Working program}
\fancyfoot[L]{\footnotesize 20 October 2026}
\fancyfoot[R]{\thepage}
\renewcommand{\headrulewidth}{0.3pt}
\setlength{\parindent}{1em}\setlength{\parskip}{4pt}
\setlength{\emergencystretch}{2em}
\setlist{nosep,leftmargin=1.5em}
\newcommand{\Z}{\mathbb Z}
\newcommand{\Q}{\mathbb Q}
\newcommand{\C}{\mathbb C}
\newcommand{\Sym}{\operatorname{Sym}}
\newcommand{\Rep}{\operatorname{Rep}}
\newcommand{\Sat}{\operatorname{Sat}}
\newcommand{\ad}{\mathrm{ad}}
\newcommand{\minrep}{\mathrm{min}}
\newcommand{\OG}{\operatorname{OG}}
\newcommand{\IG}{\operatorname{IG}}
\newcommand{\repo}{https://github.com/sean-rogers-t/Sage-programs/blob/e9fcb7f3b7d28bea9c248abab1f01d8de28250aa/}
\newcommand{\record}[2]{\href{\repo#1}{#2}}
\newcolumntype{Y}{>{\raggedright\arraybackslash}X}
\newcommand{\status}[1]{\noindent\textbf{Status:} #1\par}
\begin{document}
\begin{center}
{\LARGE\bfseries Springer--Borel maps}\\[5pt]
{\Large Integral images and the geometry of representation maps}\\[7pt]
{\normalsize A structural research program for Sean Rogers}\\[3pt]
{\small Proposal and progress inventory, 20 October 2026}
\end{center}

\section*{Purpose and central question}
The published Springer-morphism and Levi-character papers make a
representation-dependent map from characters to flag cohomology explicit.
The proposed program asks for structural control of the arithmetic and
geometry of that map:
\begin{quote}
\emph{How much of a Springer--Borel map can be predicted from the root datum,
the representation, the character lattice and the Levi subgroup, before
expanding its Schubert coefficients?}
\end{quote}
The first objective is a theory of integral images. The second studies the
generic degree and regularity of the underlying Springer morphism. These
share the same explicit map but retain different information: the Borel
quotient kills positive-degree global invariants, while the adjoint-quotient
map preserves the invariant functions needed for degree calculations.

\noindent\begin{tabularx}{\linewidth}{@{}p{0.45in}YY@{}}
\toprule
Aim & Goal & First meaningful deliverable\\
\midrule
A & Explain primes, saturation and Smith factors in integral images.
& A precise local comparison theorem and a source/parabolic criterion
for possible torsion primes.\\
B & Compute generic degrees by a uniform, effective boundary method.
& A certified face-scheme procedure, proved for a coherent family.\\
C & Determine when Springer maps preserve regular elements.
& A reusable obstruction or sufficient criterion, beyond individual examples.\\
D & Give geometric meaning to Smith filtrations.
& An integral comparison for one chain/incidence model, including primitivity.\\
E & Extend image questions to other theories and groups.
& An equivariant pilot, followed by filtered K-theory and compatible
quantum/affine constructions.\\
\bottomrule
\end{tabularx}

\paragraph{What is already available.}
The repository contains internally audited source-lattice arguments, general
normalized type B and D single-tail Smith formulas, two-step joint-image
theorems, exceptional integral-image calculations, and explicit quotient maps,
critical determinants and degree calculations for selected exceptional
representations. These are substantial starting points, with their precise
scope recorded below. They are not all published results.

\paragraph{What success means.}
A satisfactory outcome is a theorem with checkable hypotheses, a proof of the
mechanism, and independent exact examples. A counterexample delimiting a
plausible generalization, or a theorem for a well-chosen family, also counts
as progress. More tables without a question they distinguish are secondary.
This proposal organizes research; it does not commit to writing a full paper
or claim priority for the repository deductions.
Track E is now part of the program, with six dependent subprojects rather
than six simultaneous implementations.

\newpage
\section{Published foundations and common notation}
Kumar introduced the map from polynomial characters of a Levi subgroup to
flag cohomology~\cite{K}. Rogers determined the torus coordinates of the
Springer morphism by its weights~\cite{R}. Kumar--Rogers, Part II, treat
the symplectic and odd orthogonal constructions and stable injectivity;
Kumar--Xie, Part III, supply the even orthogonal analogue and exceptional
generator calculations~\cite{KR,KX}. The present program refines the
integral arithmetic and studies the geometry of the same trace map.

Let \(G\) initially be a simply connected simple complex group, with
\(T\subset B\subset G\), Weyl group \(W\), and Lie algebra
\(\mathfrak g\). Write \(X=X(T)\) for the weight lattice, \(Q_{\rm rt}\)
for the root lattice, and \(V=X\otimes\Q\). Put \(r=\operatorname{rank}G\).
We reserve \(P\) for a standard parabolic, with Levi \(L\supset T\)
and Weyl group \(W_L\). This avoids confusing the parabolic with a lattice.
Other isogeny forms are separate extensions; the orthogonal papers use their
specified \(\mathrm{SO}\) forms.

For a representation \(\rho\) with faithful differential,
\(\theta_\rho:G\to\mathfrak g\) is normalized trace projection,
so \(d\theta_\rho|_1=\mathrm{id}\). For irreducible \(V(\lambda)\)
write \(\theta_\lambda\). In the fundamental-weight/simple-coroot
coordinates of Rogers's formula,
\[
 M_\rho=\sum_\mu m_\mu\mu\mu^t,\qquad
 \theta_\rho(t)=M_\rho^{-1}\sum_\mu m_\mu\mu e^\mu(t).
\]
Here \(\mu\) is a weight column, \(m_\mu\) its multiplicity, and \(e^\mu\)
a Laurent character. This is a published coordinate method, not a new
construction of this proposal.

The pullback embeds \(\Sym_\C(\mathfrak t^*)^{W_L}\) into the character
algebra. Its image is
\(\Rep^\C_{\lambda\text{-}\mathrm{poly}}(L)\) in the papers. Their map is
\[
 \xi_\lambda^P
 =\beta^P\circ(\theta_\lambda^*)^{-1}:
 \Rep^\C_{\lambda\text{-}\mathrm{poly}}(L)
 \longrightarrow H^*(G/P,\C).
\]
The inverse is defined on the image of the injective pullback. It is not
a birational inverse of \(\theta_\lambda\).
The rational Borel map in polynomial coordinates is
\(\beta^P:\Sym_\Q(V)^{W_L}\to H^{2*}(G/P,\Q)\).
Define
\[
 \mathcal C_\rho=
 \{f\in\Sym_\Q(V):\theta_\rho^*f\in\Z[X]\},\quad
 \Lambda_\rho=(\mathcal C_\rho)_1,\quad
 \mathcal A_{\rho,L}=\mathcal C_\rho^{W_L}.
\]
Thus the literal integral character source corresponds to
\(\mathcal A_{\rho,L}\). Where integrality of \(\beta^P\) on this source
has been justified, put \(J_{\rho,L}=\beta^P(\mathcal A_{\rho,L})\).

For any specified integral image \(J_d\subset H^{2d}(G/P,\Z)\), set
\[
 S_d=\Sat(J_d)=(J_d\otimes\Q)\cap H^{2d}(G/P,\Z),\qquad
 D_d=S_d/J_d.
\]
The nonzero Smith factors describe \(D_d\), including unit factors when
recording ranks. For selected generator subalgebras, \(S_d\) need not be
all cohomology. Polynomial degree \(d\) means cohomological degree \(2d\);
it is not Laurent total degree.

\newpage
\section{Completed source work: the starting point for Aim A}
\status{The following are written, internally audited repository deductions.
Their exact literature priority and external expert review remain unresolved.}

\subsection*{2.1 The full torus source is controlled in degree one}
Under the hypotheses above, including reducible \(\rho\), the current
argument gives
\[
 \boxed{\mathcal C_\rho=\Sym_\Z(\Lambda_\rho),\qquad
 \mathcal A_{\rho,L}=\Sym_\Z(\Lambda_\rho)^{W_L}.}
\]
This includes inhomogeneous polynomials. The proof uses the full linear
integrality lattice, primitive Laurent coefficient columns and the trace
factorization to control independence modulo primes; it does not assume
that the trace form remains nondegenerate in every characteristic.
The theorem concerns characteristic-zero maps and integral coefficients,
not a universal finite-field trace-projection construction.
See \record{borel/TRACE_SOURCE_THEOREM.md}{the source theorem} and
\record{borel/TRACE_SOURCE_AUDIT.md}{its proof/source audit}
(ledger SRC-01/03).

This resolves the torus-source generation question. It does not imply that
Levi invariants are generated by invariant linear elements, or that their
Borel images are saturated. Those are the next stages of the program.

\subsection*{2.2 An irreducible highest-weight formula in every simple type}
With long roots of squared length two, let \(I_\lambda\) be the trace
index, normalized by \(I_\ad=2h^\vee\), and use the invariant inner product
to define lattice duals. The audited formula is
\[
 D_\lambda^{\rm wt}:=\sum_\mu\Z(m_\mu\mu)
       =Q_{\rm rt}+\Z\lambda,\qquad
 \boxed{\Lambda_\lambda
       =I_\lambda(Q_{\rm rt}+\Z\lambda)^*.}
\]
Multiplicity weighting means multiplying a weight by \(m_\mu\), not merely
repeating it in a list. This reduces the irreducible source calculation to
root-lattice, central-character and index data, without enumerating a large
representation. The proof uses multiplicity-one weights along suitable
lowering paths. See
\record{borel/IRREDUCIBLE_CHARACTER_LATTICES.md}{the classification} and
\record{borel/IRREDUCIBLE_LATTICE_AUDIT.md}{its audit} (SRC-02).

For \(G_2,F_4\) it specializes to \(I_\lambda Q_{\rm long}\).
For \(E_6,E_7\) it gives \(I_\lambda X\) if
\(\lambda\in Q_{\rm rt}\), and \(I_\lambda Q_{\rm rt}\) otherwise.
For \(E_8\), \(X=Q_{\rm rt}\). Here \(Q_{\rm long}\) is the lattice
generated by the long roots, in this inner-product convention.

\subsection*{2.3 Comparisons and limitations already established}
The coprime lattice/scaling lemma gives all-degree, all-Levi source and
image comparisons for the specified \(E_6/E_7\) minimal/adjoint pairs.
The \(F_4\) pair has the recorded degree-\(d\) scaling \(3^d\).
Repeating the same representation leaves normalized trace projection
unchanged. Arbitrary direct sums are subtler: three copies of the \(F_4\)
minimal module and the adjoint both have index \(18\), but their linear
sources are \(6Q_{\rm long}\) and \(18Q_{\rm long}\), respectively.
An exact combined-coefficient algorithm exists; a closed summand-level
formula remains a supporting question (SRC-04/05/07, OPEN-04/05).

\newpage
\section{Completed image work and the link to Parts II and III}
\subsection*{3.1 Classical generator formulas and integral refinements}
In the notation of the papers, type C uses \(\xi_{n,k}\), type B uses
\(\bar\xi_{n,k}\), and type D uses \(\xi^D_{n,k}\).
The targets are \(\IG(n-k,2n)\), \(\OG(n-k,2n+1)\), and
\(\OG(n-k,2n)\), respectively. The published generator and stable-injectivity
results are inputs~\cite{KR,KX}; the repository's combinatorial note gives
alternative arguments and completes the dissertation proof route in its
stated scope. It does not repair a missing proof of the published main theorems.

Our Smith results use the specified normalized tail generators
\[
 q_i=e_i(x_{n-k+1}^2,\ldots,x_n^2),\qquad
 p=x_{n-k+1}\cdots x_n,\quad p^2=q_k\quad\hbox{in type D}.
\]
They address the arithmetic of their selected image, which is additional
information beyond the complex generator formulas.

\begin{itemize}
\item \textbf{Type B:} a closed positive partition/Gaussian formula for every
nonzero Smith multiplicity, in all allowed ranks and blocks, with the
required integral primitive-prefix argument (IMG-03).
\item \textbf{Type D:} the square-sector factors agree with type B; the full
square/Euler formula includes the even-tail gluing step. The audited proof
places gluing at unit levels, with no further positive-level correction
(IMG-04/05).
\item \textbf{Two-step flags:} normalized type C joint images are saturated
and the separate-image sum is primitive. In B/D the separate sum is primitive
inside the joint image, and the ambient factors are powers of \(2\);
ambient saturation is false in general (IMG-08/09).
\end{itemize}
Here the B single-tail range has \(a=n-k\ge1,\ k\ge1\).
For D use \(n\ge4,\ a\ge1,\ k\ge2\), matching Part III.
The literal defining/vector character forms and the normalized polynomial
forms are distinguished explicitly; in the specified SO/Sp conventions the
degree-\(d\) rescaling is \(2^d\), not an implicit change of lattice.

\subsection*{3.2 Exceptional images and independent checks}
The completed finite cases include selected integral-image tables for all
\(G_2/F_4\) maximal parabolics, and the \(E_6/P_1,E_7/P_7\) cases.
The latter have primitive-generator subalgebras
\[
 \Z[h,16s]\quad\hbox{and}\quad\Z[h,81s,19683t]
\]
inside their known integral cohomology rings, with inherited relations.
Here \(h\) is the Schubert divisor; \(s\) has polynomial degree \(4\) in
\(E_6/P_1\), while \(s,t\) have polynomial degrees \(5,9\) in \(E_7/P_7\).
They are not free polynomial rings and need not be saturated.
Literal-character images and embedded minimal/adjoint comparisons are
also computed for these two spaces. Minimal and adjoint images can already
be incomparable in divisor degree.

Separately regenerated minimal-representation Springer coordinates agree
with the printed \(F_4/E_7\) formulas. The \(E_6\) comparison proposes
replacing the six denominators \(6\) by \(18\) in Part III, Lemma 7.1,
with numerators unchanged. This is a documented proposed correction,
not a published erratum. The original articles are preserved.

The \record{borel/RESEARCH_LEDGER.md}{research ledger}, entries IMG-01--09,
CASE-01--05 and LIMIT-01--05, links the exact proofs, source choices,
tables and counterexamples. None of these finite tables is being used
as an arbitrary-rank proof.

\newpage
\section{Completed geometry and the connection to Kostant--Michor}
The normalized trace construction is the generalized Cayley map considered
by Kostant--Michor~\cite{KM}. Their quotient-degree comparison and harmonic
localization theorem supply the geometric motivation. The same
\(\theta_\rho\) is being studied before passing to the Borel quotient.

\subsection*{4.1 Uniform formulas already available}
In logarithmic torus coordinates, exact differentiation gives
\[
 J_T=M_\rho^{-1}\sum_\mu m_\mu\mu\mu^t e^\mu.
\]
Equivariance determines the root-direction factors of the full,
left-translated differential:
\[
 \boxed{\left.\Psi_\rho\right|_T
 =\det(J_T)\prod_{\alpha>0}
 e^\alpha\left(\frac{\alpha(\theta_\rho)}{e^\alpha-1}\right)^2.}
\]
The root quotients extend as Laurent polynomials. Thus the Cartan Jacobian
alone is not the full determinant. A tensor-Casimir identity also provides
the quadratic invariant pullback; this applies standard representation
theory rather than asserting a new Casimir identity (GEO-02--04).

\subsection*{4.2 Exact exceptional calculations already completed}
Complete basic-invariant pullbacks and full determinant products are
packaged for \(G_2,E_6,E_7\), minimal and adjoint representations. The
\(E_6/E_7\) coordinates were checked against independent representation
matrices and exact support/unisolvence certificates. Generic degrees
on simply connected groups are recorded as follows:
\[
\begin{array}{c|rr}
 &\minrep&\ad\\ \hline
G_2&6&18\\
E_6&1080&2754\\
E_7&7272&33664
\end{array}
\]
The \(E_6/E_7\) degree arguments use normalized weight-polytope volumes,
boundary zero schemes and intersection contributions, restoring central
cover factors where appropriate. Newton volume without correction gives
the wrong answer in three of these four cases.

Seven specialized finite-field fiber reports corroborate the \(E_6/E_7\)
calculations: two primes for each of the first three cases and one for
\(E_7\) adjoint. The completed adjoint \(E_7\) root-chart fiber has length
\(16832\) at prime \(32003\); the central factor gives \(33664\).
Its staircase length was recounted independently. These reports check
specialized calculations; they are not the characteristic-zero degree proof.
See \record{borel/EXCEPTIONAL_DEGREES.md}{the degree derivation and coverage}
and the \(E_6/E_7\) geometry notes (CASE-06--08).

\subsection*{4.3 What these computations do not settle}
The source critical divisor is computed; its image is not automatically
a target branch-locus calculation, and nonproperness must be treated
separately. The degree of a particular Springer morphism is not the
minimum Cayley degree of the group. The two \(G_2\) examples have
regular-to-nonregular witnesses giving the recorded obstruction to
removing harmonic localization. Comparable general criteria remain a goal.
These distinctions define Aims B and C.

\newpage
\section{Aim A: a local theory of integral Springer--Borel images}
\status{Primary project. General source classification and several image
families are available; a uniform image criterion is not.}

\paragraph{Question.}
For the literal source \(\mathcal A_{\rho,L}\), or a separately declared
normalized generator algebra, predict the prime support of \(D_d\), give
saturation criteria, and then determine selected prime-power multiplicities.
The source choice is part of the problem, not an interchangeable convention.

\paragraph{A1. Establish a local benchmark and attribute its known content.}
Put \(K=X\cap\Lambda_\rho\) and
\(\delta_\rho=[X:K][\Lambda_\rho:K]\).
At a prime \(p\nmid\delta_\rho\), the two linear lattices agree after
localization at \(p\). The first target is a precise comparison of their Levi
invariant sources and characteristic images. Determine which additional
hypotheses, involving the characteristic map and Levi descent, imply
\[
 \beta^P\!\left(\Sym_{\Z_{(p)}}(\Lambda_\rho\otimes\Z_{(p)})^{W_L}\right)
       =H^{2*}(G/P,\Z_{(p)}).
\]
This is a \emph{target statement with hypotheses to be supplied}, not a
theorem asserted here. Classical torsion-index results and averaging at
primes not dividing \(|W_L|\) give a conservative proof route
\cite{CPZ}. First identify the part that is already a standard corollary.
The research question is how far one can remove conservative exclusions,
and what happens at the remaining primes.

\paragraph{A2. Separate three possible arithmetic obstructions.}
Compare, in fixed coordinates, the change \(X\rightsquigarrow\Lambda_\rho\),
the integral invariant-ring calculation for \(W_L\), and the Borel kernel
and integral Schubert lattice. Localize before taking a Smith form.
Use explicit invariant presentations or compatible bases to show which
obstructions survive in the image; do not assume their valuations simply add.
Work degreewise initially, and seek a ring-level criterion when possible.

\paragraph{A3. Resolve the first genuinely difficult shared-prime comparison.}
The coprime comparison lemma is already proved. Choose an exceptional
representation pair with a shared-prime obstruction and derive image
containment, sums or intersections over \(\Z_{(p)}\), accounting for the
Borel kernel. Begin on \(E_6/P_1\) or \(E_7/P_7\), where ambient rings
and independent calculations are available. Only then test additional
parabolics that distinguish competing mechanisms.

\paragraph{Proof tools and deliverables.}
Use the highest-weight lattice formula, integral invariant theory,
characteristic-map/torsion-index results, Schubert evaluation, and exact
lattice comparisons. Deliver (i) a proved local comparison with exact scope,
(ii) a finite set of possible obstruction primes under stated hypotheses,
and (iii) one sharp exceptional example with an explanation of its exponents.
A counterexample to a stronger saturation statement is a valid alternative
to (iii), provided its mechanism is proved.

\paragraph{Decision gate.}
If the first prime bound follows entirely from existing results, record
it as such and proceed to shared-prime exponents or remove an unnecessary
prime exclusion. Do not present a standard good-prime corollary as the main
new result. Related ledger questions: OPEN-02/04/05/09/10/13.

\newpage
\section{Aim B: effective generic-degree formulas from boundary geometry}
\status{The general intersection framework is available, and selected
exceptional cases are worked out; an effective uniform evaluation is missing.}

\paragraph{Question.}
Compute \(\deg\theta_\rho\) from weight and root data with an explicit account
of face degeneracies and their multiplicities. Seek a family theorem or a
certified algorithm whose output is a proof certificate, not a guessed volume.

Write \(\chi_\rho=\sum m_\mu e^\mu\) and
\(D\chi_\rho=\sum m_\mu\mu e^\mu\).
The invertible trace matrix does not affect a generic degree. On the
appropriate character lattice \(M\), the equations are
\(D\chi_\rho=b\) for generic \(b\).
Kostant--Michor identifies the full degree with the adjoint-quotient degree;
the two Weyl quotients have the same order, so no extra division by \(|W|\)
is made~\cite{KM}.

\paragraph{B1. Turn the known reduction into an effective procedure.}
For \(\Delta=\operatorname{conv}(\operatorname{supp}D\chi_\rho)\),
toric intersection theory supplies total degree
\(r!\operatorname{vol}_M(\Delta)\)~\cite{FS}.
The task is to compute all boundary contributions. On an exposed proper
face \(F\), study the scheme cut out by the face character and its
logarithmic derivatives:
\[
 \chi_F=0,\qquad D\chi_F=0.
\]
Account for the face lattice and finite covers, nonreduced structure,
positive-dimensional loci, and overlaps of strata. A symbolic
\emph{volume minus boundary} identity is a starting framework; the desired
result makes its correction terms independently evaluable.

\paragraph{B2. Prove the mechanism for a coherent family.}
Begin with minuscule representations. The \(E_6/E_7\) examples already
distinguish vanishing and nonvanishing face contributions, so
minuscule weights do not automatically mean no boundary correction.
Compare their faces with type A exterior-power weight polytopes.
Classify the singular face schemes in a restricted family before claiming
a formula for every minuscule representation. Adjoint root polytopes give
a second family once the first mechanism is understood.

\paragraph{B3. Validate the degree without conflating evidence and proof.}
Check lattice volumes by determinant triangulation, face orbits by exact
root combinatorics, and local multiplicities by an independent algebraic
description where feasible. Use exact specialized fibers, retaining
nilpotents, as corroboration. Specialization requires separate justification
if it is to enter a characteristic-zero proof. Restore a central-isogeny
degree only when the map actually factors through that cover.

\paragraph{Deliverable and fallback.}
The preferred result is a finite face-scheme/intersection procedure with
proved correctness for the selected family, plus a closed formula when its
strata admit one. A useful fallback is a complete theorem for simplex and
cross-polytope face patterns under explicit smoothness and intersection
hypotheses. \(F_4\), additional minuscule cases or \(E_8\) are tests selected
for their face geometry; an expensive \(E_8\) computation is not a milestone
by itself. Related questions: OPEN-01/11/13/16.

\newpage
\section{Aim C: regularity, critical divisors and harmonic separation}
\status{A universal determinant product and two \(G_2\) obstruction witnesses
are available. No general representation criterion is claimed.}

\paragraph{Question and consequence.}
A regular element has centralizer dimension \(\operatorname{rank}G\).
When does \(\theta_\rho\) send every regular group element to a regular
Lie-algebra element? Kostant--Michor's Theorem 3.2 identifies this condition
with their unlocalized harmonic-separation identity~\cite{KM}.
The target is a classification for a family, or a reusable obstruction,
rather than a repeated determinant expansion.

\paragraph{Route.}
For regular semisimple \(t\in T\), failure is witnessed by
\(e^\beta(t)\ne1\) for every root \(\beta\), but
\(\alpha(\theta_\rho(t))=0\) for some root \(\alpha\).
Analyze these Laurent equations and seek a root-subsystem construction
that lifts from small rank. Any lifting must prove that unwanted
root directions remain regular and that the full representation restricts
as required; a small-rank example alone does not supply this.

A semisimple witness disproves preservation. A positive theorem must also
treat regular elements with nontrivial unipotent part, using Jordan
decomposition and the published structure theory. Neither nonvanishing of
the Cartan Jacobian nor a list of full critical factors is equivalent to
that classification.

\paragraph{Deliverables.}
First recover the \(G_2\) obstruction through a formulation usable in other
types. Then prove a sufficient condition or an obstruction for a specified
representation family, using \(F_4/E_6/E_7\) as tests. Target branch and
nonproperness loci are subsequent projects, not prerequisites.
Related questions: OPEN-06/11/12.

\section{Aim D: an integral geometric explanation of Smith filtrations}
\status{Supporting project, activated when it helps explain Aim A.
The B incidence/defect interpretation is a concrete starting point.}

Hiller's chain viewpoint relates Schubert combinatorics to intersections
\cite{H}. The proposed bridge is more demanding than matching Betti numbers:
one needs a map to the \emph{integral coefficient-row filtration} of the
Springer image, with primitive prefixes and the correct powers of \(2\).

Choose one chain/incidence model and one B single-tail family. Construct
an integral comparison to the existing Schur-product prefixes, prove
primitivity geometrically, and recover the known Smith formula as a test.
Only after that attempt D's Euler gluing or B/D two-step multiplicities.
The completed even-tail theorem is not being reopened; the goal is a
conceptual explanation or a transferable proof. If only rational ranks
match, record that as insufficient, and keep the current proof.
Related questions: OPEN-03/07/08.

\newpage
\section{Track E: other cohomology theories and groups}
\status{Newly adopted research directions. The ambient theories cited below
are established; representation-dependent image extensions are proposed.
No new backend or general extension theorem is claimed.}

\subsection*{E1. Equivariant integral images: first extension to activate}
Retain the same finite-dimensional group and Springer source, but replace
ordinary cohomology by \(H_T^*(G/P,\Z)\). Put
\(S=H_T^*(\mathrm{pt},\Z)=\Sym_\Z(X)\).
The rational equivariant characteristic map
\(\beta_T^P:\Sym_\Q(V)^{W_L}\to H_T^{2*}(G/P,\Q)\)
gives a candidate \(S\)-subalgebra generated by
\(\beta_T^P(\mathcal A_{\rho,L})\), provided integrality is justified.
The equivariant parameters and polynomial Cartan inputs are different
copies of the weight variables.

The objective is to describe this embedded image as an \(S\)-module or
algebra, and explain its specialization to the ordinary image.
Positive-degree global invariant polynomials no longer simply vanish:
in the rational equivariant Borel presentation, their Cartan values agree
with the corresponding equivariant-parameter values.
Fixed-point restrictions and generalized Demazure operators provide
established computational methods~\cite{CZZ}.

\paragraph{Route and first test.}
Start with one small complete flag, then a parabolic already understood
integrally. Construct equivariant generator restrictions and a certified
module presentation. Compare exact fixed-point data with independent
equivariant Schubert multiplication. Check both the ordinary
specialization and the source normalization. Use \(G_2\) or a low-rank
classical example to distinguish mechanisms before returning to large
exceptional groups.

\paragraph{Arithmetic and the specialization problem.}
The polynomial ring \(S\) is not generally a PID, so there is no ordinary
Smith classification over \(S\). Investigate Fitting ideals, local
presentations and finite total-degree \(\Z\)-lattices. Distinguish torsion
as an \(S\)-module from torsion over \(\Z\).
The natural specialization of the image maps onto the ordinary image
when the chosen source generators specialize as intended; it need not
identify the specialized image module injectively. Kernel/base-change
phenomena must be recorded rather than inferred from a matching table.

\paragraph{Milestone and gate.}
Deliver a precisely defined integral equivariant image in one family,
its module presentation, and a proved specialization comparison.
A specialization defect with an exact explanation is an acceptable result.
Activate after Aim A's local diagram has fixed the source and integrality
conventions; do not wait for every ordinary Smith factor to be classified.
Proposed API objects should preserve the coefficient ring, parameters,
presentation and specialization maps (OPEN-14).

\paragraph{Longer objective.}
Determine whether equivariant data explains the ordinary prime support
or filtration. A claim of equivariant saturation must specify the relevant
module property; it is not an automatic translation of an ordinary
integer-lattice statement.

\newpage
\subsection*{E2. K-theory: compare character images through filtrations}
The associated-bundle map
\[
 \kappa^P:R(L)\longrightarrow K^0(G/P),\qquad
 [U]\longmapsto[G\times_P U]
\]
is defined by extending an \(L\)-representation trivially across the
unipotent radical of \(P\). Restrict it to the literal integral
Springer polynomial-character subring
\[
 R_{\rho,\mathrm{poly}}(L)
   =\theta_\rho^*(\mathcal A_{\rho,L})\subset R(L).
\]
This produces an integral image to investigate without postulating a new
Springer morphism. Its equivariant version takes values in \(K_T^0(G/P)\),
over the character ring \(R(T)=\Z[X]\). Existing equivariant K-Schubert
calculus supplies the ambient operators and classes~\cite{Uma}.

\paragraph{Question and route.}
Do the cohomological divisibility factors occur as associated-graded
features of this K-image? Compare the induced augmentation/gamma or
geometric filtrations, with the rational Chern character as a diagnostic.
The Springer polynomial grading is not automatically the augmentation
filtration, and the rational Chern character does not prove an integral
identification. State the filtration, prove its compatibility, and
measure any discrepancy.

\paragraph{Milestone.}
For one family already covered by the ordinary API, construct the
restricted K-image and compute the filtered comparison.
An exact obstruction to identifying its associated graded with
\(J_{\rho,L}\) is as valuable as a positive theorem.
Proposed library support comprises associated-bundle evaluation,
K-Demazure operations, filtered image presentations and comparison reports.
This is OPEN-17, and can start after the source conventions in Aim A
without depending on a complete equivariant implementation.

\subsection*{E3. Oriented theories: a formal-group-law framework}
An oriented theory uses
\[
 c_1(\mathcal L\otimes\mathcal M)
   =F(c_1(\mathcal L),c_1(\mathcal M)).
\]
Ordinary cohomology has \(F(u,v)=u+v\);
connective K-theory has \(F(u,v)=u+v-\beta uv\).
Algebraic cobordism supplies a universal formal group law.
Formal group algebras and formal Demazure operators already organize
characteristic maps for these theories~\cite{CPZ,CZZ}.

\paragraph{Main problem and first target.}
Define a representation-dependent source inside this framework and compare
its images under coefficient specialization. Our degree-one theorem
concerns the additive symmetric algebra; it is not already a theorem
for every formal group law. A rational logarithm may guide the comparison
but introduces denominators and does not settle integral sources.
Start with connective K-theory, explicitly specifying the coefficient
ring, completion and group-form hypotheses. Prove either a compatible
source/image specialization to the additive case or an obstruction to
such compatibility. Extend to algebraic cobordism only after that pilot
succeeds. This is the unifying longer goal behind E1/E2 (OPEN-18).

\newpage
\subsection*{E4. Quantum images: deform the map and its kernel}
The proposed target is a lift \(\xi_{\rho,q}^P\) with specialization
at \(q=0\) equal to our ordinary map. This notation denotes a goal,
not a currently defined general homomorphism. Specify its source,
coefficient ring and generator lifts before asking for its integral image.
Ordinary relations cannot be retained unchanged; for example,
\[
 H^*(\mathbb P^1,\Z)=\Z[h]/(h^2),\qquad
 QH^*(\mathbb P^1)=\Z[h,q]/(h^2-q).
\]
A map from an appropriate polynomial source can exist while its classical
kernel changes. A ring section from ordinary to quantum cohomology is
not provided by identifying their Schubert bases.

\paragraph{Route and benchmark.}
Begin with complete flags and the published quantum Borel/Toda
presentations~\cite{Kim}. Fix divisor conventions and restrict a certified
polynomial-source evaluation using quantum products. Then address a
partial-flag family, where higher-degree generator lifts need compatible
relations. Use quantum Giambelli/Pieri or Chevalley formulas as established
ambient input, not as proof that the Springer-image construction has lifted.
Track the integral normalization inherited from the original source.

\paragraph{Milestone and gate.}
Construct one multiplicative lift with a proved \(q=0\) specialization,
then compare its kernel and image with the classical ones. Study flatness
or base-change defects rather than assuming them away. Total-degree
pieces can support finite integer-lattice checks when the chosen grading
and coefficient ring make those pieces finite free. A failed candidate
lift should yield an exact relation obstruction. E4 starts after a source
and presentation are fixed, not after merely implementing quantum products
(OPEN-15).

\subsection*{E5. Affine Grassmannians: transport image questions}
The Peterson correspondence, proved by Lam--Shimozono, relates
\(QH_T^*(G/P)\) to a quotient of equivariant affine-Grassmannian homology
after specified localizations~\cite{LS}. Its product is the Pontryagin
homology product. This is not an identification with ordinary cup-product
cohomology of the affine Grassmannian.

\paragraph{Route and milestone.}
Once E4 supplies an image, transport it across the established
correspondence for a chosen \(G/P\). Record the exact quotient, inverted
classes and coefficient lattice. Describe the transported generators,
then determine what image information is lost under localization.
An integral image statement before localization requires additional
justification; a localized correspondence alone does not recover it.
The first deliverable is a commutative comparison and an exact localized
image description for one family (OPEN-19).

Separately, Mare--Mihalcea construct particular affine quantum deformations
related to periodic Toda systems~\cite{MM}. These are relevant test
frameworks, not a claim that every desired quantum theory of an arbitrary
infinite-dimensional flag variety has already been defined.

\newpage
\subsection*{E6. Kac--Moody flag varieties and extensions of the group}
There are two distinct projects. The BGG/nil-Hecke characteristic-map
side already has Kac--Moody foundations through Kostant--Kumar~\cite{KK}.
A degree-bounded backend can use nil-Hecke operators, finite Schubert
varieties or specified graded completions, with the root datum and
coefficient lattice retained. The finite-type coinvariant presentation
must not be imposed unchanged on an infinite Weyl group.

The Springer side needs a new construction. An infinite-dimensional
Kac--Moody Lie algebra cannot inject into the endomorphisms of a
finite-dimensional vector space. Highest-weight modules are typically
infinite-dimensional, so ordinary finite traces and the finite matrix
\(M_\rho=\sum m_\mu\mu\mu^t\) do not transfer directly.
An invariant bilinear form on a symmetrizable algebra by itself does
not define the desired group-to-Lie-algebra trace projection.

\paragraph{Route and first deliverable.}
First distinguish affine loop groups from general symmetrizable
Kac--Moody groups and choose a precise Cartan realization, weight lattice,
topology/completion and class of representations. Develop the
characteristic-map calculation independently of a new Springer map.
If a completed or graded-character analogue of \(\theta_\rho\) is proposed,
prove that it is well defined, equivariant and compatible with a
finite-type test or restriction. Do not call a formal infinite sum a
regular morphism. The deliverable can be a rigorous domain and
obstruction statement before any image theorem (OPEN-20).

\paragraph{Finite-dimensional group changes remain useful.}
The other-isogeny and reductive questions in OPEN-10 are closer extensions:
replace \(X(T)\) by the actual character lattice, handle central torus
directions, and state trace-nondegeneracy hypotheses.
Changing from SO to Spin is an arithmetic change of source, even when
the underlying flag varieties agree. Positive-characteristic analogues
would require separate existence, separability and trace-form analysis;
they are not supplied by reducing our characteristic-zero formulas.

\paragraph{Activation order and success criteria.}
Start E1 after Aim A fixes its local source diagram. Use E2 as a filtered
comparison pilot and E3 as the longer unifying direction.
Define E4 carefully, then use E5 as the affine bridge. E6's nil-Hecke
computations may proceed independently, but its Springer construction
is a separate foundational project.
At each stage, distinguish a known ambient theory from a new
representation-image assertion. Each pilot must yield a proved
comparison, an exact obstruction, or a precisely defined remaining gap.
None of E1--E6 is marked as implemented or proved by this proposal.

\newpage
\section{Execution plan, milestones and computational support}
The stages below are dependency gates, not elapsed-time promises.
Start with Aim A; advance B as a distinct geometric project once its next
proof obstruction is precise. C and D begin with bounded conceptual tests.
Track E follows the activation order above, starting with equivariant
images. Its later subprojects depend on their own source-definition and
specialization gates, rather than launching together.

\noindent\begin{tabularx}{\linewidth}{@{}p{0.65in}YY@{}}
\toprule
Stage & Work and output & Gate for moving on\\
\midrule
0 & Freeze notation/source forms; build theorem-to-source comparison
against Demazure/CPZ, Kumar, Hiller and Kostant--Michor. &
Each proposed claim has explicit hypotheses and a published/repository/open label.\\
1 & Prove Aim A's conservative local benchmark; compute the candidate prime
set from lattice and characteristic-map data. &
Proof separates the three arithmetic stages and identifies standard content.\\
2 & Resolve one shared-prime exceptional comparison; derive or disprove a
sharper exponent/saturation prediction. &
Independent embedded-lattice verification; theorem or exact counterexample.\\
3 & Audit and generalize the face mechanism for Aim B in a chosen family. &
All contributing strata and multiplicities accounted for; no raw-volume shortcut.\\
4 & Run the regularity-obstruction and chain-filtration pilot arguments. &
Reusable criterion/comparison, or a precise obstruction to the proposed route.\\
5 & Package established statements, certificates and limitations. &
Documentation and API expose proof scope; open questions remain separate.\\
\bottomrule
\end{tabularx}

\paragraph{First concrete work package.}
The next research session should (1) specify the local source/image diagram,
(2) identify the exact characteristic-map and Levi-descent hypotheses,
(3) prove the conservative comparison, and (4) use the existing exceptional
images to test whether each excluded prime is necessary.
The output should be a short proof note and an exact example matrix, not
a new large table campaign.

\paragraph{Existing infrastructure.}
The \texttt{sage\_borel} library already supplies Borel/BGG evaluation,
classical coinvariant normal forms, arbitrary supported standard parabolics,
rational kernels and finite-border certificates, Schubert multiplication,
inverse Borel lifts, integral image matrices and Smith forms, membership
and integer preimages, and embedded sums/intersections. Highest-weight
source lattices and the packaged exceptional geometry connect these tools
to the representation input (API-01--06).

\paragraph{Proposed API additions, conditional on the mathematics.}
Add a local source comparison report containing \(\rho\), lattice bases,
normalization, \(W_L\), degree, coefficient ring and proved prime exclusions.
Add a prime-specific image comparison returning embedded lattices and quotient
invariants, and a boundary certificate listing face lattices, scheme equations
and multiplicity justifications. A regularity witness should record the
exact torus point, root evaluations and characteristic-zero verification.
Return an explicit unresolved status when the requisite theorem is absent.
These are proposed interfaces; this document does not claim they are implemented.

\paragraph{Resources.}
Ordinary lattice and low-degree checks use existing Sage routines.
Exceptional elimination uses optional exact solvers with retained inputs,
leading ideals and logs. No external computation budget is committed.
Choose an alternative formulation or smaller discriminating case when an
elimination becomes costly; a stalled run is not a result.

\newpage
\section{Risk control, scope and evidence standards}
\paragraph{Mathematical risks and useful fallback results.}
Integral invariant rings can obscure a simple source classification.
Fallback: a prime-local theorem or one Levi family.
Boundary schemes may be nonisolated or overlap. Fallback: a family with
proved transversality/local models, rather than dropping those strata.
Regularity preservation may fail broadly. Fallback: a general obstruction
with exact witnesses, rather than seeking a positive theorem at all costs.
A chain model may recover ranks without integral primitivity. Fallback:
identify the missing comparison explicitly and retain the existing proof.

\paragraph{Verification and literature.}
Every general theorem needs a written argument with its integral steps
exposed. Verify decisive examples through a different realization:
character/coefficient calculations versus representation matrices;
coinvariant arithmetic versus independent cohomology presentations;
geometric degree arguments versus exact fiber calculations.
Agreement between implementations sharing the same input is a weaker check.
Before a novelty claim, compare the exact statement and hypotheses with the
relevant literature; an unsuccessful search does not establish originality.

The current baseline is repository commit
\texttt{e9fcb7f} (20 October 2026), not an assertion about what has merged
to the default branch. The runtime is version \(0.20.0\).
The existing release record binds 157 regression checks and 24 isolated-wheel
checks to runtime hashes; documentation verification records 40 executable
examples and 908 local links/anchors. These establish computational coverage,
not proof validity, external review or research priority.

\paragraph{Tracking progress.}
Use the \record{borel/RESEARCH_LEDGER.md}{research ledger} and
\record{borel/PROOF_STATUS.md}{detailed proof-status record}.
For each milestone record: exact question; hypotheses; proof status;
independent checks; source/novelty comparison; remaining obstruction;
next discriminating test. Mark a theorem complete only when its general
argument is supplied. A finite example stays a finite example.
The proposal itself adds no mathematical theorem.

\noindent\begin{tabularx}{\linewidth}{@{}p{1.4in}Y@{}}
\toprule
Ledger questions & Role in this program\\
\midrule
OPEN-02/05/09 & Core Aim A: arithmetic, shared primes and Levi images.\\
OPEN-04/10 & Source extensions supporting A: direct sums and isogeny/reductive scope.\\
OPEN-01 & Core Aim B: effective generic-degree calculation.\\
OPEN-06/12 & Core Aim C: critical divisors and regularity/localization.\\
OPEN-03 & Core Aim D: integral chain/incidence explanation.\\
OPEN-07/08 & Conditional next applications: B/D joint factors and many Levi blocks.\\
OPEN-11/13 & Target geometry and further exceptional tests, chosen by a mechanism.\\
OPEN-14/17/18 & Track E1--E3: equivariant, filtered K-theoretic
and formal-group-law images.\\
OPEN-15/19/20 & Track E4--E6: quantum lifts, affine transport and
Kac--Moody characteristic/Springer constructions.\\
OPEN-16 & Deferred: comparisons with minimum Cayley degrees.\\
\bottomrule
\end{tabularx}

\paragraph{Expected contribution.}
The intended advance is a passage from explicit maps to a local arithmetic
theory of their images, together with an effective geometric degree method.
The completed source classification provides the representation input;
the classical and exceptional examples provide demanding tests.
A successful program explains which arithmetic and geometric features are
forced, which depend on the representation or parabolic, and where a proposed
general rule fails.

\newpage
\section*{Sources and reproducible starting material}
The published bibliography below supplies the construction and background.
Repository notes are unpublished working documents, with internal audits.
Their presence here does not establish priority. The original article PDFs
and their provenance are indexed in the
\record{articles/README.md}{article collection}.
The detailed literature tasks at Stage 0 remain necessary; this is a
source-grounded proposal, not an exhaustive literature review.

\begin{thebibliography}{99}\small\raggedright
\setlength{\itemsep}{1pt}
\bibitem{K}
S.~Kumar, \emph{Representation ring of Levi subgroups versus cohomology ring
of flag varieties}, Math.\ Ann.\ \textbf{366} (2016), 395--415.
\href{https://doi.org/10.1007/s00208-015-1325-6}{DOI};
\href{https://kumar.math.unc.edu/papers/kumar76.pdf}{author PDF}.

\bibitem{R}
S.~Rogers, \emph{An explicit determination of the Springer morphism},
Comm.\ Algebra \textbf{46} (2018), 4233--4242.
\href{https://doi.org/10.1080/00927872.2017.1407418}{DOI};
\href{https://arxiv.org/abs/1701.01538}{arXiv:1701.01538}.
The repository copy consulted is the 2017 preprint, not the final journal PDF.

\bibitem{KR}
S.~Kumar and S.~Rogers, \emph{Representation ring of Levi subgroups versus
cohomology ring of flag varieties II}, J.\ Algebra \textbf{556} (2020), 340--362.
\href{https://doi.org/10.1016/j.jalgebra.2020.02.029}{DOI};
\href{https://arxiv.org/abs/1907.10089}{arXiv:1907.10089}.

\bibitem{KX}
S.~Kumar and J.~Xie, \emph{Representation ring of Levi subgroups versus
cohomology ring of flag varieties III}, J.\ Pure Appl.\ Algebra
\textbf{228} (2024), 107539.
\href{https://doi.org/10.1016/j.jpaa.2023.107539}{DOI};
\href{https://kumar.math.unc.edu/papers/kumar86.pdf}{author PDF}.

\bibitem{KM}
B.~Kostant and P.~W.~Michor, \emph{The generalized Cayley map from an
algebraic group to its Lie algebra}, in \emph{The Orbit Method in Geometry
and Physics}, Progress in Mathematics \textbf{213} (2003), 259--296.
\href{https://arxiv.org/abs/math/0109066}{arXiv:math/0109066};
especially Theorem 3.2 and Corollary 3.3.

\bibitem{CPZ}
B.~Calm\`es, V.~Petrov and K.~Zainoulline,
\emph{Invariants, torsion indices and oriented cohomology of complete flags},
Ann.\ Sci.\ \'{E}c.\ Norm.\ Sup\'er.\ (4) \textbf{46} (2013), no.~3.
\href{https://arxiv.org/abs/0905.1341}{arXiv:0905.1341}.
Includes the relation to Demazure's characteristic-map theory;
all coefficient hypotheses must be retained when applying it.

\bibitem{H}
H.~Hiller, \emph{Combinatorics and intersections of Schubert varieties},
Comment.\ Math.\ Helv.\ \textbf{57} (1982), 41--59.
\href{https://www.e-periodica.ch/cntmng?pid=com-001:1982:57::10}{journal archive}.
A chain/intersection background source; an integral Smith comparison
is a proposed task, not a theorem attributed to this paper.

\bibitem{FS}
W.~Fulton and B.~Sturmfels, \emph{Intersection theory on toric varieties},
\href{https://arxiv.org/abs/alg-geom/9403002}{arXiv:alg-geom/9403002}.
Toric intersection background; the representation-specific boundary
analysis is recorded separately.

\bibitem{CZZ}
B.~Calm\`es, K.~Zainoulline and C.~Zhong,
\emph{Equivariant oriented cohomology of flag varieties},
Documenta Math., Extra Volume: Alexander S.~Merkurjev's Sixtieth Birthday
(2015), 113--144.
\href{https://arxiv.org/abs/1409.7111}{arXiv:1409.7111}.

\bibitem{Uma}
V.~Uma, \emph{Equivariant K-theory of flag varieties revisited and
related results}, Colloq.\ Math.\ \textbf{132} (2013), 151--175.
\href{https://arxiv.org/abs/1101.5911}{arXiv:1101.5911}.

\bibitem{Kim}
B.~Kim, \emph{Quantum cohomology of flag manifolds \(G/B\) and quantum
Toda lattices}, Ann.\ of Math.\ (2) \textbf{149} (1999), 129--148.
\href{https://arxiv.org/abs/alg-geom/9607001}{arXiv:alg-geom/9607001}.

\bibitem{LS}
T.~Lam and M.~Shimozono, \emph{Quantum cohomology of \(G/P\) and homology
of affine Grassmannian}.
\href{https://arxiv.org/abs/0705.1386}{arXiv:0705.1386}.

\bibitem{MM}
A.-L.~Mare and L.~C.~Mihalcea,
\emph{An affine deformation of the quantum cohomology ring of flag
manifolds and periodic Toda lattice}.
\href{https://arxiv.org/abs/1409.3587}{arXiv:1409.3587};
\href{https://doi.org/10.1112/plms.12077}{DOI}.

\bibitem{KK}
B.~Kostant and S.~Kumar, \emph{The nil Hecke ring and cohomology of
\(G/P\) for a Kac--Moody group \(G\)}, Proc.\ Natl.\ Acad.\ Sci.\ USA
\textbf{83} (1986), 1543--1545.
\href{https://doi.org/10.1073/pnas.83.6.1543}{DOI};
\href{https://pubmed.ncbi.nlm.nih.gov/16593661/}{article record}.
\end{thebibliography}

\noindent\textbf{Working-note map at the baseline commit.}
\begin{itemize}\small
\item Source arguments: \record{borel/TRACE_SOURCE_THEOREM.md}{degree-one theorem},
\record{borel/TRACE_SOURCE_AUDIT.md}{source audit},
\record{borel/IRREDUCIBLE_CHARACTER_LATTICES.md}{irreducible lattice formula},
\record{borel/IRREDUCIBLE_LATTICE_AUDIT.md}{lattice audit}.
\item Classical arithmetic:
\record{borel/TYPE_B_GENERAL_SMITH.md}{B theorem},
\record{borel/TYPE_D_SQUARE_THEOREM.md}{D square comparison},
\record{borel/TYPE_D_EVEN_GLUING_THEOREM.md}{D full gluing theorem},
\record{borel/TYPE_C_JOINT_PRIMITIVITY.md}{C joint images},
\record{borel/ORTHOGONAL_JOINT_IMAGES.md}{B/D joint images}.
\item Exceptional arithmetic:
\record{borel/EXCEPTIONAL_MINUSCULE.md}{selected minuscule images},
\record{borel/E6_CHARACTER_FORM.md}{E6 character form},
\record{borel/E7_CHARACTER_FORM.md}{E7 character form}.
\item Geometry and computation:
\record{borel/SPRINGER_GEOMETRY.md}{general/G2 geometry},
\record{borel/E6_GEOMETRY.md}{E6 geometry},
\record{borel/E7_GEOMETRY.md}{E7 geometry},
\record{borel/EXCEPTIONAL_DEGREES.md}{degrees and verification coverage},
\record{borel/release-0.20-verification.json}{release audit}.
\end{itemize}
\end{document}
